% !TEX root = ../Hauptdatei.tex
% ============================================================
% Kapitel: Inline- und Display-Modus
% ============================================================

\chapter{Inline- und Display-Modus}
\label{sec:modi}
\index{Mathematikmodus}

\section{Inline-Mathematik}
\index{Inline-Mathematik}

Mathematische Ausdrücke im Fließtext werden mit \texttt{\$...\$}\index{\$ (Inline-Mathematik)} gesetzt:
Die Gleichung $E = mc^2$ ist wohl die bekannteste Formel der Physik.
Der Satz des Pythagoras lautet $a^2 + b^2 = c^2$.

\section{Display-Modus}
\index{Display-Modus}

Abgesetzte Formeln mit \texttt{\textbackslash[...\textbackslash]}:
\[
    E = mc^2
\]

Nummerierte Gleichungen mit der \texttt{equation}-Umgebung\index{equation (Umgebung)} \cite{mittelbach2004}:
\begin{equation}
    a^2 + b^2 = c^2
    \label{eq:pythagoras}
\end{equation}

% ============================================================
\chapter{Hoch- und Tiefstellungen}
\label{sec:stellen}
\index{Hochstellung}\index{Tiefstellung}
% ============================================================

\section{Exponenten und Indizes}

\begin{itemize}
    \item Exponent: $x^2$, $x^{10}$, $e^{i\pi}$
    \item Index: $a_1$, $a_{n+1}$, $x_{\max}$
    \item Beides: $x_i^2$, $a_{i,j}^{(k)}$
\end{itemize}

\section{Summen, Produkte und Integrale}

\[
    \sum_{i=1}^{n} i = \frac{n(n+1)}{2}
\]

\[
    \prod_{k=1}^{n} k = n!
\]

\[
    \int_{0}^{\infty} e^{-x^2}\, dx = \frac{\sqrt{\pi}}{2}
\]

\section{Grenzwerte}
\index{Grenzwert}

\[
    \lim_{n \to \infty} \left(1 + \frac{1}{n}\right)^n = e
\]

% ============================================================
\chapter{Brüche und Wurzeln}
\label{sec:brueche}
\index{Bruch}\index{Wurzel}
% ============================================================

\section{Verschiedene Brucharten}

\begin{equation}
    \frac{a}{b} \qquad \dfrac{a}{b} \qquad \tfrac{a}{b}
    \label{eq:brucharten}
\end{equation}

\begin{itemize}
    \item \texttt{\textbackslash frac}\index{frac} -- Standardbruch
    \item \texttt{\textbackslash dfrac}\index{dfrac} -- Großer Bruch (Display-Stil)
    \item \texttt{\textbackslash tfrac}\index{tfrac} -- Kleiner Bruch (Text-Stil)
\end{itemize}

\section{Geschachtelte Brüche}

\begin{equation}
    1 + \frac{1}{1 + \frac{1}{1 + \frac{1}{1 + \cdots}}}
    \label{eq:kettenbruch}
\end{equation}

\section{Wurzeln}

\[
    \sqrt{2} \approx 1{,}414 \qquad
    \sqrt[3]{8} = 2 \qquad
    \sqrt[n]{a} = a^{1/n}
\]

% ============================================================
\chapter{Matrizen}
\label{sec:matrizen}
\index{Matrix}
% ============================================================

\section{Einfache Matrix}

\[
    \begin{matrix}
        1 & 2 & 3 \\
        4 & 5 & 6 \\
        7 & 8 & 9
    \end{matrix}
\]

\section{Matrix mit Klammern}

Verschiedene Klammerarten:

\[
    \begin{pmatrix} 1 & 0 \\ 0 & 1 \end{pmatrix}
    \quad
    \begin{bmatrix} 1 & 0 \\ 0 & 1 \end{bmatrix}
    \quad
    \begin{vmatrix} a & b \\ c & d \end{vmatrix}
    \quad
    \begin{Vmatrix} 1 & 0 \\ 0 & 1 \end{Vmatrix}
\]

\begin{itemize}
    \item \texttt{pmatrix}\index{pmatrix} -- runde Klammern
    \item \texttt{bmatrix}\index{bmatrix} -- eckige Klammern
    \item \texttt{vmatrix}\index{vmatrix} -- einfache Striche (Determinante)
    \item \texttt{Vmatrix}\index{Vmatrix} -- doppelte Striche
\end{itemize}

\section{Große Matrix mit Punkten}

\[
    A = \begin{pmatrix}
        a_{11} & a_{12} & \cdots & a_{1n} \\
        a_{21} & a_{22} & \cdots & a_{2n} \\
        \vdots & \vdots & \ddots & \vdots \\
        a_{n1} & a_{n2} & \cdots & a_{nn}
    \end{pmatrix}
\]

% ============================================================
\chapter{Mehrzeilige Gleichungen}
\label{sec:mehrzeilig}
\index{Gleichung!mehrzeilig}
% ============================================================

\section{align -- ausgerichtete Gleichungen}

\begin{align}
    (a + b)^2 &= a^2 + 2ab + b^2 \label{eq:binom1} \\
    (a - b)^2 &= a^2 - 2ab + b^2 \label{eq:binom2} \\
    a^2 - b^2 &= (a+b)(a-b) \label{eq:binom3}
\end{align}

\section{alignat -- mehrfache Ausrichtung}

\begin{alignat*}{2}
    &\text{Addition:} &\quad a + b &= c \\
    &\text{Subtraktion:} &\quad a - b &= d \\
    &\text{Multiplikation:} &\quad a \cdot b &= e
\end{alignat*}

\section{cases -- Fallunterscheidungen}

\begin{equation}
    |x| = \begin{cases}
        x & \text{falls } x \geq 0 \\
        -x & \text{falls } x < 0
    \end{cases}
    \label{eq:absolut}
\end{equation}

\section{split -- lange Gleichung}

\begin{equation}
    \begin{split}
        (a+b)^3 &= a^3 + 3a^2b + 3ab^2 + b^3 \\
                 &= a^3 + b^3 + 3ab(a+b)
    \end{split}
    \label{eq:binom3d}
\end{equation}

% ============================================================
\chapter{Griechische Buchstaben und Symbole}
\label{sec:symbole}
\index{Griechische Buchstaben}\index{Symbol}
% ============================================================

\section{Griechische Buchstaben}

\begin{align*}
    &\alpha \quad \beta \quad \gamma \quad \delta \quad \epsilon \quad \varepsilon \\
    &\zeta \quad \eta \quad \theta \quad \vartheta \quad \iota \quad \kappa \\
    &\lambda \quad \mu \quad \nu \quad \xi \quad \pi \quad \varpi \\
    &\rho \quad \varrho \quad \sigma \quad \varsigma \quad \tau \quad \upsilon \\
    &\phi \quad \varphi \quad \chi \quad \psi \quad \omega \\
    &\Gamma \quad \Delta \quad \Theta \quad \Lambda \quad \Xi \quad \Pi \\
    &\Sigma \quad \Upsilon \quad \Phi \quad \Psi \quad \Omega
\end{align*}

\section{Häufige Operatoren und Symbole}

\[
    \forall \quad \exists \quad \nexists \quad \in \quad \notin \quad \subset \quad \subseteq
\]
\[
    \cup \quad \cap \quad \setminus \quad \emptyset \quad \mathbb{N} \quad \mathbb{Z} \quad \mathbb{R}
\]
\[
    \Rightarrow \quad \Leftrightarrow \quad \iff \quad \implies \quad \approx \quad \equiv \quad \neq
\]
\[
    \leq \quad \geq \quad \ll \quad \gg \quad \propto \quad \perp \quad \parallel
\]

\section{Pfeile}

\[
    \leftarrow \quad \rightarrow \quad \Leftarrow \quad \Rightarrow \quad
    \leftrightarrow \quad \Leftrightarrow
\]
\[
    \mapsto \quad \longmapsto \quad \hookrightarrow \quad \uparrow \quad \downarrow
\]

% ============================================================
\chapter{Auszüge und Theoreme}
\label{sec:theoreme}
\index{Theorem}\index{Definition}\index{Beweis}
% ============================================================

\section{Mathematische Umgebungen}

\begin{definition}
Eine Funktion $f\colon \mathbb{R} \to \mathbb{R}$ heißt \emph{stetig}
an der Stelle $x_0$, wenn gilt:
\[
    \lim_{x \to x_0} f(x) = f(x_0)
\]
\end{definition}

\begin{satz}[Satz von Pythagoras]
In einem rechtwinkligen Dreieck mit den Katheten $a$ und $b$ und der
Hypotenuse $c$ gilt:
\begin{equation}
    a^2 + b^2 = c^2
\end{equation}
\end{satz}

\section{Beweisumgebung}

\begin{beweis}
Sei $ABC$ ein rechtwinkliges Dreieck mit dem rechten Winkel bei $C$.
Dann gilt nach dem Kathetensatz die Behauptung.
\end{beweis}

% ============================================================
\chapter{Referenzen auf Gleichungen}
\label{sec:math-referenzen}
\index{Referenz!Gleichung}
% ============================================================

Wie in den vorherigen Abschnitten gezeigt, lassen sich Gleichungen
referenzieren:

\begin{itemize}
    \item Gleichung~\ref{eq:pythagoras}: Der Satz des Pythagoras
    \item Gleichung~\ref{eq:kettenbruch}: Ein Kettenbruch
    \item Gleichungen~\ref{eq:binom1}--\ref{eq:binom3}: Binomische Formeln
    \item Gleichung~\ref{eq:absolut}: Die Betragsfunktion
\end{itemize}